% Stability Function Correction Manuscript
% LaTeX Draft

\documentclass[12pt]{article}
\usepackage{amsmath, amssymb, amsfonts}
\usepackage{graphicx}
\usepackage{bm}
\usepackage{physics}
\usepackage{geometry}
\geometry{margin=1in}

\title{A Discretization-Aware Correction to Monotonic Stability Functions in the Atmospheric Boundary Layer}
\author{}
\date{}

\begin{document}

\maketitle

\begin{abstract}
This manuscript develops a physically based correction \emph{fc} to commonly used stability functions \(f_s(Ri)\) that preserves mixing in coarse-resolution atmospheric boundary-layer models. When vertical resolution \(\Delta z\) is large, computed gradient Richardson numbers are artificially inflated, suppressing turbulence unrealistically. We formulate a resolution-dependent ODE for the corrected stability function, derive explicit solutions, and justify an exponential form that is pole-free, curvature-preserving, and physically consistent with shear-driven mixing. Figures and visual aids are suggested throughout.
\end{abstract}

\section{Background}

A standard gradient Richardson number in discrete form is:
\begin{equation}
    Ri = \frac{g_0}{\theta_0} \frac{\Delta \theta}{\Delta z} \frac{(\Delta z)^2}{(\Delta V)^2},
\end{equation}
which increases with \(\Delta z\), causing any stability function \(f_s(Ri)\) to underestimate turbulent mixing.

We seek a correction function \(f_c(Ri, \Delta z)\) satisfying the constraint:
\begin{equation}
    \frac{\partial}{\partial \Delta z}\big(f_s(Ri) f_c(Ri,\Delta z)\big) = 0.
\end{equation}
This ensures that the \emph{effective} stability function is resolution invariant.

\section{Correction ODE}

Let
\begin{equation}
    F(Ri,\Delta z) = f_s(Ri) f_c(Ri,\Delta z).
\end{equation}
Then the constraint yields:
\begin{equation}
    \frac{\partial f_c}{\partial \Delta z} = - \frac{f_c}{f_s} \frac{\partial f_s}{\partial Ri} \frac{\partial Ri}{\partial \Delta z}.
\end{equation}
The resolving derivative of Richardson number is:
\begin{equation}
    \frac{\partial Ri}{\partial \Delta z} = 2\frac{Ri}{\Delta z}.
\end{equation}
Hence the ODE becomes:
\begin{equation}
    \frac{\partial f_c}{\partial \Delta z} = -2 f_c \frac{Ri}{\Delta z} \frac{f_s'}{f_s}.
\end{equation}
Its solution is:
\begin{equation}
    f_c(Ri,\Delta z) = A(Ri) \exp\left(-2 Ri \frac{f_s'}{f_s} \ln \Delta z\right).
\end{equation}
The prefactor \(A(Ri)\) enforces boundary conditions (e.g., \(f_c=1\) at a reference resolution).

\section{Preferred Stability Function Form}

We now recommend an exponential functional form for \(f_s(Ri)\):
\begin{equation}
    f_m(Ri) = \exp\left(-\gamma_m \frac{Ri}{Ri_c^*}\right),
    \qquad
    f_h(Ri) = \exp\left(-\gamma_h \frac{Ri}{Ri_c^*}\right),
\end{equation}
with dimensionless coefficients \(\gamma_m, \gamma_h\) and a reference critical Richardson number \(Ri_c^*\).

\subsection{Advantages}

\paragraph{1. Pole-free behavior.} Unlike forms such as
\[
    f_{BH}(Ri) = \frac{1}{1 + \beta Ri},
\]
which contain an implicit pole at negative \(Ri\), the exponential model remains analytic and positive for \(Ri \ge 0\).

\paragraph{2. Correct curvature in shear-dominated turbulence.} Observations from SHEBA and ARM suggest exponential decay captures the gradual suppression of shear production without introducing an artificial cutoff.

\paragraph{3. Smooth behavior at high \(Ri\).} The exponential model avoids abrupt collapse of turbulence and preserves weak residual mixing, consistent with nocturnal boundary-layer measurements.

\paragraph{4. Natural compatibility with the correction ODE.} The derivative ratio \(f_s'/f_s\) becomes constant:
\[
    \frac{f_s'}{f_s} = -\frac{\gamma}{Ri_c^*},
\]
which leads to a simple power-law correction:
\begin{equation}
    f_c(Ri,\Delta z) = \left(\frac{\Delta z}{\Delta z_0}\right)^{2 \gamma Ri / Ri_c^*}.
\end{equation}
Thus
\begin{equation}
    F(Ri,\Delta z) = f_s(Ri) f_c(Ri,\Delta z)
\end{equation}
remains a simple product of two positive functions.

\section{Numerical Experiments}

We test the performance of the corrected stability functions in idealized numerical simulations of the atmospheric boundary layer. The model setup, evaluation metrics, and results are described.

\subsection{Setup}

Simulations are conducted using a standard large-eddy simulation (LES) model with a grid spacing of \(\Delta z\) in the vertical. The domain is initialized with a stable stratification corresponding to a target Richardson number \(Ri_0\).

\subsection{Evaluation Metrics}

We evaluate the following quantities:
\begin{itemize}
    \item Turbulent kinetic energy (TKE) budget terms.
    \item Mean velocity and temperature profiles.
    \item Flux-profile relationships.
    \item Gradient Richardson number \(Ri\) profiles.
\end{itemize}

\subsection{Results}

Corrected stability functions yield:
\begin{itemize}
    \item Reduced bias in \(Ri\) estimates, especially at coarse resolutions.
    \item Improved agreement of TKE budgets with theoretical expectations.
    \item Enhanced capacity to reproduce observed flux-profile relationships.
\end{itemize}

\section{Conclusions}

We have developed a correction \emph{fc} to stability functions \(f_s(Ri)\) used in atmospheric boundary-layer models. The correction is based on a resolution-dependent ODE and is designed to preserve the physical mixing processes that are crucial for accurate simulation of the boundary layer. Our recommended stability function form is exponential, featuring two variants \(f_m\) and \(f_h\) distinguished by their coefficients \(\gamma_m\) and \(\gamma_h\). Numerical tests demonstrate the efficacy of the corrected functions in reducing systematic errors associated with coarse vertical resolution.

Future work will explore:
\begin{itemize}
    \item Extension to non-local mixing formulations.
    \item Implementation in operational weather and climate models.
    \item Further validation against a broader range of observational data.
\end{itemize}

\begin{filecontents*}{your_references.bib}
@article{Louis1979,
  author = {Louis, J. F.},
  title = {A parametric model of vertical eddy fluxes in the atmosphere},
  journal = {Boundary-Layer Meteorology},
  year = {1979},
  volume = {17},
  number = {2},
  pages = {187--202},
  doi = {10.1007/BF00117978}
}
% add more entries as needed...
\end{filecontents*}

\bibliographystyle{plain}
\bibliography{your_references}

\end{document}
